\section{Conclusions}\label{sec:concl}
We proposed HiPhyBoost, a hierarchical, physics-anchored boosting algorithm for coseismic landslide forecasting, and tested it against physics-based, global-style logistic and gradient-boosting models on nine Pacific Rim earthquakes, including the 2016 Kaik\={o}ura event, using only open and rapidly available Earth-observation inputs. The main conclusions are as follows.
\begin{enumerate}
\item Transferred models rank hillslopes reasonably well, but they mispredict the overall level of landsliding in a new earthquake by up to three orders of magnitude. Calibration slopes close to one show that the error is a shift in level, not in shape.
\item Separating the shared, physics-anchored response from latent event offsets improves transfer. HiPhyBoost raises the median leave-one-event-out AUC from 0.805 (physics-constrained boosting) to 0.833 and reaches AUC 0.917 for Kaik\={o}ura without using any New Zealand data. It is better in six of nine events, but with nine events the gain is not statistically significant ($p=0.16$) and remains to be confirmed on further inventories.
\item The event offsets summarise the landslide propensity of an earthquake relative to generic Newmark mechanics and ShakeMap shaking. They range over 5.7 log-odds units, from Denali (lowest) to Cinchona (highest). Their order is consistent with what is known about antecedent wetness, weak volcanic soils and ShakeMap bias, but they also absorb differences in mapping completeness and study-domain definition, which nine events cannot separate.
\item Before any landslide is mapped, the empirical-Bayes prior over offsets gives a predictive interval for the landslide total that contained the observed total in seven of nine events. Baseline models give only a point estimate, which was off by a factor of 20--50 in the worst cases.
\item Rapid mapping of a few image tiles corrects much of the level error, but only if the spatial clustering of landslides is respected. Tile-level landslide counts are overdispersed by a factor of 7--190 relative to independent cells. A cell-level Bayesian update is therefore badly overconfident: its 90\% intervals contained the truth in at most 26\% of trials. HiPhyBoost's quasi-posterior keeps the best point accuracy of all methods (median error a factor of 1.6 after one tile and 1.4 after ten) and raises coverage to 0.72--0.76, while its interval for the landslide total narrows from a factor of about 60 to about 4 between one and ten mapped tiles.
\item Monotonic physical constraints prevent physically implausible responses to shaking, which occurred in a quarter to almost half of held-out cells for unconstrained models. They do not make SHAP attributions more stable across events: the form of the shaking response transfers, but its location on the shaking axis does not.
\end{enumerate}
For Aotearoa New Zealand, the results support using Pacific Rim training data for scenario forecasting of future Alpine Fault or Hikurangi earthquakes, provided that forecasts are issued as intervals and that rapid mapping of a few image tiles is built into the response plan to narrow them.
